\chapter{Resultats i Avaluació de les Heurístiques}
\label{chap:resultats}

Aquest capítol presenta els resultats obtinguts de l'avaluació exhaustiva d'ambdós entorns de simulació (Singles i Doubles). S'analitza l'efectivitat i fiabilitat de les versions estructurades respecte als baselines, utilitzant les eines de \textit{benchmark} internes de l'entorn Core, desenvolupades sota el disseny \textit{Round-Robin} automatitzat explicat prèviament.

\section{Avaluació Exhaustiva al Format 1 vs 1 (Singles)}

L'anàlisi sistemàtic del format \textit{Singles} va procedir a confrontar iterativament tots els agents heurístics dissenyats (des de \textit{v1} fins a \textit{v6}) i els \textit{baselines} absoluts de l'entorn.

Les proves es van estructurar en sèries massives, llançant matrius d'enfrontaments per cada possible par (ex. $v6$ vs \textit{Random}, $v3$ vs \textit{MaxBasePower}), on els milers d'execucions es van processar i sintetitzar de forma paral·lela obtenint el guanyador o perdedor teòric. En l'observació del percentatge de victòries net (\textit{Win Rate}), la tendència generalitzadora marca clarament una \textbf{progressió evolutiva de disseny encertada}. 
Generacions inicials com la $v1$ no obtenien victòries aclaparadores més enllà del 55\%-60\% en intentar basar decisions únicament en regles ofendents de dany simple respecte atacs a l'atzar del \textit{RandomPlayer}. La introducció en modelat teòric profund amb \textit{The Tracker} ($v3$) i l'establiment expert d'incorporació de tipologies defensives del motor al \textit{The Champion} ($v6$) augmentava de cop en un rendiment tàcit els índexs per sobre del 90\%-95\% quant es combatien davant els elements de \textit{poke-env}.

\textbf{El pes de la Estocasticitat}: Malgrat crear intel·ligències heurístiques altament sofisticades, de la visió general numèrica se n'extreu que una \textbf{victòria del 100\% és intrínsecament impossible en aquest entorn i tipus d'agents reaccionaris purs}. Això resideix essencialment a la incommensurable potència de la inobservabilitat, afegit al RNG absolut del joc i al dany latent d'encerts crítics aleatoris que es poden atorgar de manera favorable totalment a un element de sort inferior que utilitzi combinatòries subòptimes (\textit{baselines}). L'anàlisi de matrius presentades indica clarament aquesta resistència asimptòtica.

\vspace{0.5cm}
\begin{figure}[htbp]
    \centering
    % Espai preparat per afegir la matriu i la gràfica v3 del repertori 1vs1
    \includegraphics[width=0.85\linewidth, draft]{benchmarks_singles_v3/placeholder_singles_heatmap.png}
    \caption{Matriu de calor mostrant el percentatge de victòries (\textit{Win Rate}) de tots els encreuaments de \textit{Singles} evaluant baselines vs. $v1$-$v6$. Es pot lligar com els iteratius majors assoleixen superioritat al marge inferior esquerre.}
    \label{fig:singles_results}
\end{figure}

\section{Avaluació al Format 2 vs 2 (Doubles)}

De manera similar, es va testejar la fiabilitat al \textit{Core Doubles} aprofitant el disseny \textit{Score-then-Combine}. Sota la multiplicitat d'arbre de l'enginyeria, el càlcul d'eficiència es restringia exclusivament en els prototips primigenis del llinatge \textit{v1} al \textit{v3} per simplificació experimental i focalització d'estudi previ al model DRL per xarxes neurals.

La visualització empírica als tornejos demostrà de nou que hi hagué una fita creixent en els \textit{baselines} quan es presentaven múltiples entitats objectives (\textit{multi-target}). Exigir decisions conjuntes penalitza radicalment als \textit{RandomPlayers} on sovint aquests malgastaven els atacs atacant-se simultàniament les pròpies files si es produïen errors del llançament a l'atzar respecte el tipus objectiu (foc amic).
No obstant això, el rendiment estabilitzador a través de combinant parelles òptimes per pes amb el \textit{v3 Doubles} garanteix ser base per establir estratègies pre-predictores suficients i emmagatzemar la pre-avaluació requerida que encoratjarà bases de model entrenables i més complexes.   

\vspace{0.5cm}
\begin{figure}[htbp]
    \centering
    % Espai preparat per afegir la visualització procedent de s02_doubles/results/ o data/benchmarks_doubles_v2
    \includegraphics[width=0.85\linewidth, draft]{benchmarks_doubles_v2/placeholder_doubles_heatmap.png}
    \caption{Avaluació exhaustiva al format \textit{Doubles}. La progressió de $v1$ fins a $v3$ en el comportament de \textit{Score-then-Combine} contra els \textit{baselines} teòrics de \textit{poke-env}.}
    \label{fig:doubles_results}
\end{figure}
